% =====================================================================
\chapter{Simuladores diferenciables y modelos sustitutos}
\label{tema:3}
% =====================================================================

El \cref{tema:2} obtuvo el gradiente de un simulador de elementos finitos derivando a mano sus ecuaciones. Este tema generaliza la idea a cualquier simulador, estacionario o temporal, y presenta la alternativa cuando derivar no es posible o el simulador es demasiado caro: aprender un modelo sustituto con pocas simulaciones y optimizar sobre él. Se siguen las notas de S.~G. Johnson sobre adjuntos para recurrencias~\cite{johnson}, las de cálculo matricial del MIT~\cite[§6.3]{mit18s096}, los capítulos de modelos sustitutos de Kochenderfer y Wheeler~\cite{kw}, y los artículos de DiffTaichi~\cite{difftaichi}, las redes informadas por la física~\cite{pinn}, el operador neuronal de Fourier~\cite{fno} y MeshGraphNets~\cite{mgn}. La física de la aleta procede del capítulo~3 de Incropera \emph{et al.}~\cite{incropera}.

\section{Derivar a través de un simulador}

Casi todo simulador discreto resuelve un sistema de ecuaciones $R(u,\theta)=0$, donde $u\in\R^m$ es el estado (temperaturas, desplazamientos, velocidades) y $\theta\in\R^n$ son los parámetros de diseño.

\begin{teorema}[Sensibilidad por la función implícita y el adjunto]\label{teo:implicita}
Sean $R:\R^m\times\R^n\to\R^m$ y $J(\theta)=g\bigl(u(\theta),\theta\bigr)$ funciones continuamente derivables, con $R(u(\theta),\theta)=0$ y $\partial R/\partial u$ invertible en la solución. Entonces
\[
\frac{du}{d\theta}=-\Bigl(\frac{\partial R}{\partial u}\Bigr)^{-1}\frac{\partial R}{\partial\theta},\qquad
\frac{dJ}{d\theta}=\frac{\partial g}{\partial\theta}-\lambda^\T\frac{\partial R}{\partial\theta},\quad\text{con}\quad\Bigl(\frac{\partial R}{\partial u}\Bigr)^{\!\T}\lambda=\Bigl(\frac{\partial g}{\partial u}\Bigr)^{\!\T}.
\]
\end{teorema}

\begin{proof}
El teorema de la función implícita garantiza que $u(\theta)$ es derivable cerca de la solución. Derivando la identidad $R(u(\theta),\theta)=0$ se obtiene $\frac{\partial R}{\partial u}\frac{du}{d\theta}+\frac{\partial R}{\partial\theta}=0$, de donde sale la primera fórmula. Sustituyéndola en $\frac{dJ}{d\theta}=\frac{\partial g}{\partial\theta}+\frac{\partial g}{\partial u}\frac{du}{d\theta}$ y agrupando $\frac{\partial g}{\partial u}\bigl(\frac{\partial R}{\partial u}\bigr)^{-1}=\lambda^\T$ se obtiene la segunda.
\end{proof}

El \cref{teo:adjunto} es el caso $R=K(x)U-F$. El coste del adjunto es una resolución con la matriz transpuesta, que en problemas lineales puede reutilizar la factorización ya calculada; el de la sensibilidad directa, $n$ resoluciones. Para simuladores temporales, $u_{k+1}=\Phi(u_k,\theta)$, el mismo razonamiento da una recurrencia adjunta que se recorre hacia atrás en el tiempo~\cite{johnson}:
\[
\lambda_N=\Bigl(\frac{\partial g}{\partial u_N}\Bigr)^{\!\T},\qquad \lambda_{k-1}=\Bigl(\frac{\partial\Phi}{\partial u}\Bigr)^{\!\T}_{\!k}\lambda_k,\qquad \frac{dg}{d\theta}=\frac{\partial g}{\partial\theta}+\sum_{k=1}^N\lambda_k^\T\frac{\partial\Phi}{\partial\theta}\Big|_{u_{k-1}} .
\]
Es la retropropagación a través del tiempo de una red recurrente (módulo~12, §5.3), con los mismos problemas: hay que guardar todos los estados, lo que se alivia con el recálculo por puntos de control (módulo~10, §5.2), y en dinámicas caóticas o con muchos pasos el gradiente puede explotar.

\section{Diferenciación automática de simuladores}

En lugar de derivar a mano, el simulador se puede escribir en un lenguaje que propague derivadas. JAX transforma programas de Python y NumPy (derivación, vectorización y compilación para GPU y TPU)~\cite{jax}; NVIDIA Warp compila núcleos de Python para GPU que son derivables~\cite{warp}; Taichi, con su extensión DiffTaichi, permite escribir simuladores físicos derivables, y sus autores informan de que su simulador elástico necesita unas 4,2 veces menos líneas que una versión CUDA escrita a mano, con velocidad similar, y es 188 veces más rápido que la implementación en TensorFlow~\cite{difftaichi}. PyTorch también sirve cuando el simulador se expresa con operaciones de tensores, como en el caso de estudio de este tema.

Tres advertencias limitan el enfoque. Los eventos discretos (contactos, cambios de malla, condicionales sobre el estado) producen funciones discontinuas cuya derivada automática es correcta pero inútil, porque vale cero o no ve el salto. Las derivadas de un resolvedor iterativo no deben calcularse a través de sus iteraciones, sino con el \cref{teo:implicita} en la solución convergida. Y el coste de memoria del modo inverso crece con el número de pasos de tiempo.

\section{Modelos sustitutos}

\begin{definicion}[Modelo sustituto]\label{def:sustituto}
Sea $s:\Theta\to\R$ una respuesta cara de evaluar (el resultado de un simulador). Un \emph{modelo sustituto} es una función $\hat s_\varphi$ barata, ajustada con un conjunto de diseño $\{(\theta_i,s(\theta_i))\}_{i=1}^M$, con la que se sustituye a $s$ dentro de la optimización.
\end{definicion}

Los sustitutos clásicos son las superficies de respuesta polinómicas ajustadas por mínimos cuadrados y los procesos gaussianos, que además estiman su incertidumbre~\cite{kw}. Los neuronales cubren desde un perceptrón que aproxima una respuesta escalar hasta modelos que aproximan campos completos: las redes informadas por la física (PINN) minimizan el residuo de la ecuación diferencial en puntos de colocación, calculando las derivadas de la red por diferenciación automática~\cite{pinn}; los operadores neuronales de Fourier aprenden la aplicación entre funciones de entrada y de salida, y sus autores informan de aceleraciones de hasta tres órdenes de magnitud frente a resolvedores clásicos~\cite{fno}; MeshGraphNets aprende la dinámica sobre la malla de simulación con una red de grafos~\cite{mgn}. Bibliotecas como NVIDIA PhysicsNeMo (antes Modulus) agrupan estas arquitecturas para construir sustitutos de CFD, mecánica estructural o análisis térmico~\cite{physicsnemo}.

El diseño de experimentos importa tanto como el modelo. El muestreo por hipercubo latino divide el rango de cada variable en $M$ intervalos y toma exactamente una muestra por intervalo, de modo que sus proyecciones sobre cada eje quedan uniformemente repartidas~\cite{kw}.

\begin{proposicion}[Error del óptimo de un sustituto]\label{prop:2eps}
Sea $\Theta$ compacto, $s$ y $\hat s$ continuas y $\sup_{\theta\in\Theta}\abs{\hat s(\theta)-s(\theta)}\le\varepsilon$. Si $\theta^\star$ minimiza $s$ y $\hat\theta$ minimiza $\hat s$ en $\Theta$, entonces
\[
s(\theta^\star)\le s(\hat\theta)\le s(\theta^\star)+2\varepsilon .
\]
\end{proposicion}

\begin{proof}
La primera desigualdad es la optimalidad de $\theta^\star$. Para la segunda, $s(\hat\theta)\le\hat s(\hat\theta)+\varepsilon\le\hat s(\theta^\star)+\varepsilon\le s(\theta^\star)+2\varepsilon$, usando sucesivamente la cota de error, la optimalidad de $\hat\theta$ para $\hat s$ y otra vez la cota.
\end{proof}

La proposición es tranquilizadora y peligrosa a la vez: la garantía exige un error \emph{uniforme} en todo $\Theta$, mientras que la validación solo mide el error en unos puntos. El optimizador busca precisamente las zonas en que $\hat s$ es optimista, que suelen ser aquellas en que hay pocos datos. Es el mismo fenómeno que la sobreoptimización de un modelo de recompensa (módulo~21, §6.4). Por eso todo diseño propuesto por un sustituto se verifica con el simulador y, si hay presupuesto, se añade a los datos y se reajusta el modelo.

\section{Caso de estudio: el perfil de una aleta disipadora}

\begin{ejemplo}[Aleta de aluminio con sustituto neuronal y con simulador diferenciable]\label{ej:aleta}
\emph{Contexto.} Un disipador de aluminio ($k=200$ W/(m·K)) refrigerado por aire forzado ($h=60$ W/(m²·K)) debe evacuar calor desde una base a $\theta_b=T_b-T_\infty=50$ K. La masa de cada aleta está fijada: por cada metro de profundidad, el área del perfil es $A=60$ mm². El perfil es trapecial, con espesor $t(x)=t_b(1-s\,x/L)$ para $0\le x\le L$, donde $L$ es la longitud y $s\in[0,1)$ la conicidad; la condición de área fija da $t_b=A/\bigl(L(1-s/2)\bigr)$. Por extrusión, la punta debe medir al menos $t_{\min}=0{,}3$ mm.

\emph{Modelo físico.} Con aleta delgada (número de Biot $ht/(2k)\ll 1$), conducción longitudinal y convección por las dos caras, la temperatura $\theta(x)$ por metro de profundidad satisface~\cite{incropera}
\[
\frac{d}{dx}\Bigl(k\,t(x)\,\frac{d\theta}{dx}\Bigr)-2h\,\theta=0,\qquad \theta(0)=\theta_b,\qquad \theta'(L)=0,
\]
y el calor evacuado es $q=2h\int_0^L\theta\,dx$. Para $s=0$ hay solución cerrada, $q=\sqrt{2hkt}\,\theta_b\tanh(mL)$ con $m=\sqrt{2h/(kt)}$.

\emph{Formulación.}
\[
\max_{L,\,s}\ q(L,s)\quad\text{sujeto a}\quad t_b(1-s)\ge t_{\min},\quad 20\le L\le 100\ \text{mm},\quad 0\le s\le 0{,}9 .
\]
El simulador discretiza la ecuación por volúmenes finitos con 400 celdas y resuelve un sistema tridiagonal; reproduce la solución cerrada de la aleta rectangular con un error relativo de $1{,}4\cdot 10^{-6}$.

\emph{Línea base.} Para $s=0$, el óptimo con área fija cumple $\sinh(2\beta)=6\beta$ con $\beta=mL$ (\cref{eje:3-schmidt}), es decir, $\beta=1{,}4192$. Con los datos, $L=58{,}6$ mm, $t=1{,}02$ mm y $q=220{,}4$ W/m: unos 11 W para una aleta de 50 mm de profundidad.

\emph{Método 1: sustituto neuronal.} Se supone que cada simulación es cara (en un diseño real sería un cálculo CFD conjugado de horas) y se dispone de 40. Se toman 60 puntos por hipercubo latino en $(L,s)$, se descartan los no fabricables y se simulan los 40 primeros factibles. Un perceptrón 2--32--32--1 con activación tanh, entrenado con 32 puntos, tiene un error cuadrático medio de 10,1 W/m (4,9~\% de la media) en los 8 reservados. Reentrenado con los 40, su máximo en una rejilla factible de $201\times 201$ está en $L=62{,}4$ mm y $s=0{,}815$, con predicción de 242,2 W/m; la simulación de verificación (la número 41) da 242,0 W/m.

\emph{Método 2: simulador diferenciable.} El mismo simulador escrito con tensores de PyTorch (200 celdas, sistema denso) da $\partial q/\partial L$ y $\partial q/\partial s$ por diferenciación automática, que coinciden con diferencias finitas centradas en seis cifras. Ascenso de gradiente con Adam (módulo~7, §5.3) y \emph{proyección} sobre la restricción de la punta, $s\le s_{\max}(L)=(1-r)/(1-r/2)$ con $r=t_{\min}L/A$, converge en 150 iteraciones a $L=66{,}3$ mm y $s=0{,}801$, con $q=242{,}2$ W/m y la restricción activa (punta de 0,300 mm).

\emph{Comprobación y conclusión.} Una búsqueda exhaustiva en $81\times 96$ puntos (7776 simulaciones) da $q=242{,}1$ W/m en $L=66$ mm y $s=0{,}80$. Los dos métodos llegan a menos del 0,1~\% de ese valor y mejoran la aleta rectangular óptima en un 9,8~\% con la misma masa. El óptimo está sobre la restricción de fabricación: si se pudiera hacer la punta más fina, la aleta mejoraría hacia el perfil triangular. El sustituto encontró una solución algo distinta ($L=62$ mm) con casi el mismo calor porque la función es muy plana cerca del óptimo; el método del gradiente necesita el simulador derivable, pero no un presupuesto de muestras.
\end{ejemplo}

\begin{table}[h]
\centering\small
\begin{tabularx}{\textwidth}{@{}l>{\raggedright\arraybackslash}X>{\raggedright\arraybackslash}p{2.3cm}>{\raggedright\arraybackslash}X@{}}
\toprule
Herramienta & Para qué sirve & Licencia & Cuándo elegirla \\
\midrule
JAX & derivación, vectorización y compilación de código NumPy & Apache-2.0 & simuladores nuevos escritos en estilo funcional \\
NVIDIA Warp & núcleos de simulación derivables en GPU & Apache-2.0 & física de partículas, telas o sólidos en GPU \\
Taichi (DiffTaichi) & lenguaje para simulación física derivable & Apache-2.0 & simuladores de partículas y rejillas de alto rendimiento \\
PyTorch & tensores y diferenciación automática & BSD-3 (con componentes de otras licencias) & simuladores pequeños unidos a redes neuronales \\
NVIDIA PhysicsNeMo & PINN, operadores neuronales y GNN para física & Apache-2.0 & sustitutos de CFD o de estructuras a escala industrial \\
scikit-learn & MLP, procesos gaussianos, validación & BSD-3 & sustitutos de pocas variables con pocos datos \\
dolfin-adjoint & adjuntos automáticos de modelos FEniCS & LGPL & EDP generales con elementos finitos \\
\bottomrule
\end{tabularx}
\caption{Tecnologías del tema~3. Licencias según los metadatos publicados por cada proyecto.}
\label{tab:tec3}
\end{table}

\begin{observacion}[Errores frecuentes en la práctica]\label{obs:errores3}
(i) \emph{Sobreajuste del sustituto}: con 40 puntos y miles de pesos, la red interpola el ruido; hay que reservar datos de validación y regularizar. (ii) \emph{Extrapolación}: el optimizador lleva el diseño a los bordes del dominio, donde hay menos datos; se verifica siempre con el simulador (\cref{prop:2eps}). (iii) \emph{Fugas de datos}: si los puntos de validación son casi idénticos a los de entrenamiento (simulaciones repetidas con otra malla, por ejemplo), el error de validación es optimista. (iv) \emph{La verdad tampoco es exacta}: el simulador tiene error de discretización; hay que comprobar la convergencia de malla antes de entrenar con él. (v) \emph{Validez del modelo físico}: el caso supone aleta delgada; con $ht/k$ grande o radiación apreciable, el modelo unidimensional no vale. (vi) \emph{Restricciones de seguridad}: un sustituto puede predecir un diseño factible que no lo es; las restricciones críticas se comprueban con el simulador o con un margen. (vii) \emph{Coste mal contado}: el entrenamiento del sustituto y las simulaciones de verificación forman parte del presupuesto.
\end{observacion}

\begin{lecturas}
\item M.~J. Kochenderfer y T.~A. Wheeler~\cite{kw}: caps. «Sampling Plans», «Surrogate Models», «Probabilistic Surrogate Models» y «Surrogate Optimization».
\item S.~G. Johnson~\cite{johnson}: adjuntos para relaciones de recurrencia (MIT 18.335J).
\item Y.~Hu \emph{et al.}~\cite{difftaichi}: programación diferenciable para simulación física.
\item M.~Raissi, P.~Perdikaris y G.~E. Karniadakis~\cite{pinn}; Z.~Li \emph{et al.}~\cite{fno}; T.~Pfaff \emph{et al.}~\cite{mgn}.
\item F.~P. Incropera \emph{et al.}~\cite{incropera}, cap.~3: conducción unidimensional y superficies extendidas.
\end{lecturas}

\section*{Ejercicios del tema 3}
\addcontentsline{toc}{section}{Ejercicios del tema 3}

\begin{ejercicio}{1}{eje:3-elegir}
Para cada caso, elija entre sustituto, simulador diferenciable o diferencias finitas y justifíquelo: (a) 3 parámetros y un CFD comercial de cajas negras que tarda 6 horas; (b) 5000 densidades y un código de elementos finitos propio; (c) un simulador de choques con contactos y 20 parámetros.
\end{ejercicio}

\begin{ejercicio}{1}{eje:3-pinn}
Escriba la pérdida de una red informada por la física $\theta_\varphi(x)$ para la aleta rectangular del \cref{ej:aleta}: identifique el término de residuo y los de condiciones de contorno, y explique por qué no hace falta malla.
\end{ejercicio}

\begin{ejercicio}{2}{eje:3-aleta}
Para una aleta rectangular de $L=40$ mm y $t=1{,}5$ mm ($A=60$ mm²), con los datos del \cref{ej:aleta}, calcule $m$, $mL$, el calor $q$, la eficiencia $\eta=\tanh(mL)/(mL)$ y el número de Biot $ht/(2k)$.
\fuente{adaptado de F.~P. Incropera \emph{et al.}, \emph{Fundamentals of Heat and Mass Transfer}, cap.~3 (aletas de sección uniforme)~\cite{incropera}.}
\end{ejercicio}

\begin{ejercicio}{2}{eje:3-coste}
La viga MBB tiene 1200 variables y cada resolución de $KU=F$ tarda 10 ms. Compare el tiempo de calcular el gradiente por diferencias finitas hacia delante y por el método adjunto, y el de 100 iteraciones de optimización con cada uno.
\end{ejercicio}

\begin{ejercicio}{2}{eje:3-implicita}
Sea $R(u,\theta)=A(\theta)u-b$ con $A(\theta)=\begin{psmallmatrix}2+\theta&-1\\-1&2\end{psmallmatrix}$, $b=(1,0)^\T$ y $J=u_1+u_2$. En $\theta=0$, calcule $u$, el vector adjunto $\lambda$ y $dJ/d\theta$, y compruébelo con una diferencia finita de paso $0{,}01$.
\end{ejercicio}

\begin{ejercicio}{2}{eje:3-cuadratico}
Se simula la aleta rectangular del \cref{ej:aleta} en $L=30$, 60 y 90 mm, con resultados $q=165{,}38$; $220{,}34$ y $198{,}20$ W/m. Ajuste un sustituto cuadrático $\hat q(L)=\beta_0+\beta_1L+\beta_2L^2$ ($L$ en mm), calcule su máximo y compárelo con el óptimo verdadero ($L=58{,}6$ mm, $q=220{,}42$ W/m), sabiendo que el simulador da $218{,}13$ W/m en el máximo del sustituto.
\fuente{adaptado de M.~J. Kochenderfer y T.~A. Wheeler, \emph{Algorithms for Optimization}, cap. «Surrogate Models»~\cite{kw}.}
\end{ejercicio}

\begin{ejercicio}{2}{eje:3-recurrencia}
Un componente se enfría según $\theta_{k+1}=(1-a\Delta t)\,\theta_k$ con $\theta_0=50$ K, $\Delta t=1$ s y $a=0{,}1$ s$^{-1}$. Calcule $\theta_3$ y $d\theta_3/da$ derivando la expresión cerrada y con la recurrencia adjunta hacia atrás.
\fuente{adaptado de S.~G. Johnson, «Adjoint methods and sensitivity analysis for recurrence relations», notas de MIT 18.335J (otoño de 2010)~\cite{johnson}.}
\end{ejercicio}

\begin{ejercicio}{3}{eje:3-cota}
(a) Demuestre la \cref{prop:2eps}. (b) Construya un ejemplo con $\Theta=[0,1]$ en el que el sustituto tenga error nulo en tres puntos de validación y, sin embargo, su minimizador sea pésimo para $s$. ¿Qué enseña sobre la validación de sustitutos?
\end{ejercicio}

\begin{ejercicio}{3}{eje:3-schmidt}
(a) Demuestre que, para la aleta rectangular de área de perfil fija $A$, el calor $q(L)=\sqrt{2hkA/L}\,\theta_b\tanh\bigl(c\,L^{3/2}\bigr)$, con $c=\sqrt{2h/(kA)}$, es máximo cuando $\beta=cL^{3/2}$ cumple $\sinh(2\beta)=6\beta$. (b) Con $h=40$ W/(m²·K), $k=170$ W/(m·K), $A=50$ mm² y $\theta_b=50$ K, calcule $L$, $t$ y $q$ óptimos, usando $\beta=1{,}4192$.
\end{ejercicio}
